\section{面積の変化率}\label{節:導入基本定理}
ここまでの2つの問い,接線と面積は,一見すると別の話である.接線では縦の変化を横の変化で割り,面積では縦の長さと横の長さを掛けて足している.この2つに関係はあるのだろうか.面積の右端を固定せず,動かしてみよう.

$y=x^2$ の下で,\,$0$ から $x$ までの面積を $A(x)$ と書く.ただし,この段階では $x\ge0$ とする.右端を $x$ から $x+h$ へ少し動かすと,面積の増加 $A(x+h)-A(x)$ は,幅 $h$ の細い帯の面積である.\,$h>0$ なら,その帯の高さは $x^2$ 以上 $(x+h)^2$ 以下だから,
\[
 x^2h\le A(x+h)-A(x)\le(x+h)^2h.
\]
正の数 $h$ で割ると,
\[
 x^2\le\frac{A(x+h)-A(x)}h\le(x+h)^2.
\]
$h$ を $0$ に近づけると左右はともに $x^2$ へ近づく.\footnote{$x>0$ で $-x\le h<0$ のときは,帯の幅が $-h$,高さが $(x+h)^2$ 以上 $x^2$ 以下なので $(x+h)^2(-h)\le A(x)-A(x+h)\le x^2(-h)$ となる.正の数 $-h$ で割れば $(x+h)^2\le\frac{A(x+h)-A(x)}h\le x^2$ であり,左側から近づけても同じく $x^2$ へ近づく.}面積の平均変化率も,その間にあるので $x^2$ へ近づくはずである.つまり,面積の瞬間の増加率は,その場所の曲線の高さになるのだ.

そうであれば,微分して $x^2$ になる関数を見つけることで,長方形の和を計算せずに面積の候補が得られるだろう.例えば $F(x)=\frac{x^3}{3}$ を候補として調べてみよう.平均変化率は
\begin{align*}
 \frac{F(x+h)-F(x)}h
 &=\frac{(x+h)^3-x^3}{3h}\\
 &=\frac{3x^2h+3xh^2+h^3}{3h}\\
 &=x^2+xh+\frac{h^2}{3}
\end{align*}
であり,確かに $h$ を $0$ へ近づけると $x^2$ になる.さらに $F(0)=0$ であり,\,$A(0)=0$ と同じ値をとる.\,$x=1$ とすれば $F(1)=\frac13$ となり,長方形と三角形で求めた値に一致する.ただし,この計算は $F$ が「瞬間の増加率が $x^2$」「$x=0$ で $0$」という2つの条件に合うことを示しているだけで,その条件で面積が1つに決まることはまだ証明していない.後で平均値の定理を学ぶと,「微分して同じ関数になる2つの関数の差は定数である」と示せる.それを使って,\,$A(x)=F(x)$ であると結論できる.

ここでは,候補を見つける計算と,候補が正しいと確かめる論理を分けておこう.積分の計算で「微分の逆を探す」という方針を使えるのは,この両方が成り立つからである.
